\startSection
\selectlanguage{italian}
\sectionTitle{L'Inno Nazionale Italiano}

\begin{question}{L'inno nazionale italiano è conosciuto anche come \enquote{Fratelli d'Italia}. Chi scrisse il testo?}
  \choice{Giuseppe Verdi}
  \choice{Alessandro Manzoni}
  \choice[correct]{Goffredo Mameli}
  \choice{Giuseppe Garibaldi}
  \choice{Gabriele D'Annunzio}
\end{question}

\begin{question}{Chi compose la musica dell'inno conosciuto come \enquote{Il Canto degli Italiani}?}
  \choice{Gioachino Rossini}
  \choice[correct]{Michele Novaro}
  \choice{Giuseppe Verdi}
  \choice{Vincenzo Bellini}
  \choice{Pietro Mascagni}
\end{question}

\begin{question}{In quale anno Goffredo Mameli compose il testo dell'inno?}
  \choice{1846}
  \choice[correct]{1847}
  \choice{1848}
  \choice{1861}
  \choice{1870}
\end{question}

\begin{question}{Quale tra i seguenti versi apre ufficialmente l'inno nazionale italiano?}
  \choice{\enquote{Italia, Italia, sacra terra}}
  \choice{\enquote{Viva l'Italia unita e forte}}
  \choice[correct]{\enquote{Fratelli d'Italia, l'Italia s'è desta}}
  \choice{\enquote{Gloria all'Italia nostra}}
  \choice{\enquote{Sorgi Italia, dal tuo sonno}}
\end{question}

\begin{question}{In quale contesto storico nacque il testo di Mameli?}
  \choice{La ricostruzione dopo la Seconda guerra mondiale}
  \choice{La fondazione dell'Impero romano}
  \choice[correct]{Il Risorgimento e le lotte per l'unità nazionale}
  \choice{La firma dei Patti Lateranensi}
  \choice{La nascita dell'Unione Europea}
\end{question}

\begin{question}{Perché l'inno di Mameli fu riconosciuto definitivamente come inno nazionale solo in epoca recente?}
  \choice{Perché era rimasto sconosciuto fino al XXI secolo.}
  \choice{Perché non aveva una melodia ufficiale.}
  \choice{Perché era scritto in latino classico.}
  \choice[correct]{Perché per lungo tempo fu usato in via provvisoria prima della piena formalizzazione.}
  \choice{Perché era riservato esclusivamente alle cerimonie sportive.}
\end{question}
\shuffleSection
\renderSection